\documentclass[11pt]{article}
\usepackage[ngerman]{babel}
\usepackage[a4paper,margin=1.0in]{geometry}
\usepackage[T1]{fontenc}
\usepackage[utf8]{inputenc}
\usepackage{lmodern}
\usepackage{amsmath,amssymb,amsthm,mathrsfs}
\usepackage{microtype}
\usepackage{fancyhdr}
\usepackage{array,longtable,enumitem}
\setlength{\parindent}{1.5em}
\setlength{\parskip}{0.18em}
\emergencystretch=3em
\pagestyle{fancy}
\fancyhf{}
\cfoot{\thepage}
\lhead{Heinrich Weber, Lehrbuch der Algebra, Band II}
\rhead{wortgetreue Reparaturabschrift}
\newcommand{\Om}{\Omega}
\newcommand{\Omm}{\Omega_m}
\newcommand{\Hm}{H_m}
\newcommand{\eps}{\varepsilon}
\newcommand{\srcpage}[1]{\par\medskip\noindent\textsf{[Quellseite~#1]}\par\smallskip}
\begin{document}

\section*{\S\,200. Normaleinheiten in $H_m$.}

\srcpage{728}
Wir führen jetzt eine vereinfachende Bezeichnung ein, indem wir
\begin{equation}
r^{5^j}=r_j
\tag{1}
\end{equation}
setzen, und $\beta$ ein volles Restsystem nach dem Modul $\nu$ durchlaufen lassen. Es folgt hieraus
\begin{equation}
r_{\beta+\nu'}=-r_\beta,
\tag{2}
\end{equation}
und in der Form $(r,r_j)$ sind dann zwar nicht alle Substitutionen der Gruppe des Körpers $\Omega_m$ enthalten, sondern es müssen noch die Substitutionen $(r,r_j^{-1})$ hinzukommen; wohl aber liefert uns $(r,r_j)$ alle Substitutionen der Gruppe von $H_m$.

Wir bezeichnen jetzt mit $\mathfrak E(r)$ eine Einheit des Körpers $H_m$. Unter diesen sind auch die Einheiten des Körpers $H_\mu$ (als
\srcpage{729}
imprimitive Zahlen) enthalten und durch die Gleichung $\mathfrak E(r)=\mathfrak E(-r)$ charakterisirt.

Wir wollen nun unter einer \textbf{Normaleinheit} des Körpers $H_m$ eine Einheit verstehen, die nicht $=\pm1$ ist, aber der Gleichung
\begin{equation}
\mathfrak E(r)\mathfrak E(-r)=\pm1
\tag{3}
\end{equation}
genügt.

Es ist klar, dass eine Einheit des Körpers $H_\mu$ dieser Bedingung jedenfalls nicht genügen kann. Aber nicht jede Einheit in $H_m$, die dem Körper $H_\mu$ nicht angehört, wird Normaleinheit sein\footnote{Der Verfasser hat in der oben erwähnten Arbeit (Acta mathematica, Bd. 8) die Normaleinheiten als „primitive Einheiten“ bezeichnet. Dieser Name ist hier verworfen, weil er zu dem Irrthum Veranlassung geben könnte, als ob jede Einheit des Körpers $H_m$, die nicht zugleich dem Körper $H_\mu$ angehört, dazu zu rechnen sei.}.

Die Einheiten $\tau_j$ [\S~199, (5)] gehören aber zu den Normaleinheiten. Es ist nämlich
\begin{equation}
\tau_j=\tg\frac{5^j\pi}{m}=\frac{1-r_j}{1+r_j}r_j,
\tag{4}
\end{equation}
und wenn darin die Substitution $(r,-r)=(r_\beta,r_{\beta+\nu'})$ gemacht wird, so geht $\tau_j$ in
\begin{equation}
\tau_{\beta+\nu'}=-\frac1{\tau_\beta}
\tag{5}
\end{equation}
über, was der Relation (3) entspricht.

Ein System von $\nu'$ Normaleinheiten
\[
\mathfrak E_0(r),\ \mathfrak E_1(r),\ldots,\mathfrak E_{\nu'-1}(r)
\]
soll ein unabhängiges System heissen, wenn die Determinante
\begin{equation}
\pm\sum\pm\log|\mathfrak E_0(r_0)|\log|\mathfrak E_1(r_1)|\cdots\log|\mathfrak E_{\nu'-1}(r_{\nu'-1})|
=L(\mathfrak E_0,\mathfrak E_1,\ldots,\mathfrak E_{\nu'-1})
\tag{6}
\end{equation}
einen von Null verschiedenen Werth hat.

Den absoluten Werth dieser Determinante, $L$, nennen wir auch hier den \textbf{Regulator} des Systems $\mathfrak E_0,\mathfrak E_1,\ldots,\mathfrak E_{\nu'-1}$.

Die Einheiten $\tau_0,\tau_1,\ldots,\tau_{\nu'-1}$ bilden ein unabhängiges System normaler Einheiten, denn ihr Regulator
\[
L(\tau_0,\tau_1,\ldots,\tau_{\nu'-1})
\]
ist eben die im vorigen Paragraphen definirte Determinante $T$ [\S~199, (14), (15)]. Und dass diese nicht Null sein kann, ergiebt sich unmittelbar daraus, dass sie in dem Ausdrucke für die
\srcpage{730}
Classenzahl als Factor vorkommt. Die Classenzahl ist aber, da gewiss immer wenigstens eine Classe, die Hauptclasse, existirt, von Null verschieden, und also kann auch $T$ nicht gleich Null sein. Wir haben daher den Satz:

\begin{center}\textbf{1. Die Zahlen $\tau_0,\tau_1,\ldots,\tau_{\nu'-1}$ sind ein unabhängiges System von Normaleinheiten des Körpers $H_m$.}\end{center}

Ist $\mathfrak E(r)$ eine Normaleinheit, so ist, wenn wir
\begin{equation}
l_\beta=\log|\mathfrak E(r_\beta)|
\tag{7}
\end{equation}
setzen, wegen (3)
\begin{equation}
l_{\beta+\nu'}=-l_\beta.
\tag{8}
\end{equation}

Ist nun $\mathfrak E_0,\mathfrak E_1,\ldots,\mathfrak E_{\nu'-1}$ irgend ein System unabhängiger Normaleinheiten, so setzen wir
\begin{equation}
l_{i\beta}=\log|\mathfrak E_i(r_\beta)|,\qquad l_{i,\beta+\nu'}=-l_{i\beta},
\tag{9}
\end{equation}
so dass der Regulator des Systems der $\mathfrak E_i$ den Ausdruck erhält:
\begin{equation}
\pm\sum\pm l_{00}l_{11}\cdots l_{\nu'-1,\nu'-1}=L(\mathfrak E_0,\mathfrak E_1,\ldots,\mathfrak E_{\nu'-1})
\tag{10}
\end{equation}
und eine von Null verschiedene Zahl ist, die wir mit $T_\mathfrak E$ bezeichnen.

Bedeutet $\mathfrak E(r)$ irgend eine andere Normaleinheit, so können wir die Unbekannten $\xi_0,\xi_1,\ldots,\xi_{\nu'-1}$ aus den $\nu'$ linearen Gleichungen bestimmen, die wir erhalten, wenn wir in
\begin{equation}
l_\beta=\xi_0l_{0\beta}+\xi_1l_{1\beta}+\cdots+\xi_{\nu'-1}l_{\nu'-1,\beta}
\tag{11}
\end{equation}
den Index $\beta$ von $0$ bis $\nu'-1$ gehen lassen, und wegen (8) und (9) ist diese Gleichung auch für die übrigen Werthe von $\beta$ richtig, so dass man daraus die für alle $r$ gültige Formel ableiten kann:
\begin{equation}
\mathfrak E(r)=\pm\mathfrak E_0(r)^{\xi_0}\mathfrak E_1(r)^{\xi_1}\cdots\mathfrak E_{\nu'-1}(r)^{\xi_{\nu'-1}}.
\tag{12}
\end{equation}

Bedeuten andererseits $m_0,m_1,\\ldots,m_{\nu'-1}$ irgend welche ganze rationale Zahlen, so sind alle in der Form
\[
\pm\mathfrak E_0(r)^{m_0}\mathfrak E_1(r)^{m_1}\cdots\mathfrak E_{\nu'-1}(r)^{m_{\nu'-1}}
\]
enthaltenen Zahlen Normaleinheiten des Körpers $H_m$. Demnach lassen sich auf das Gleichungssystem (10) die Betrachtungen anwenden, die wir im \S~186 ausführlich durchgeführt haben, woraus sich ergiebt:

\begin{center}\textbf{2. Die aus (10) oder (11) bestimmten Exponenten $\xi_0,\xi_1,\ldots,\xi_{\nu'-1}$ sind rationale Zahlen.}\end{center}

\srcpage{731}
Hieraus kann, wie im \S~187, gefolgert werden, dass es Systeme von $\nu'$ unabhängigen Normaleinheiten giebt, deren Regulator $T_0$ so klein als möglich wird, und solche Systeme heissen \textbf{Fundamentalsysteme von Normaleinheiten}.

Ist $\mathfrak E_0,\mathfrak E_1,\ldots,\mathfrak E_{\nu'-1}$ ein solches Fundamentalsystem, so ergeben sich die Exponenten $\xi_0,\xi_1,\ldots,\xi_{\nu'-1}$ in (10) und (11) als \textbf{ganze rationale Zahlen}. Auch dies kann ebenso bewiesen werden, wie im \S~187.

Stellen wir unter dieser Voraussetzung über die $\mathfrak E_i$ die Gleichungen (11) für irgend ein anderes System unabhängiger Einheiten $\mathfrak E_0^{(1)},\mathfrak E_1^{(1)},\ldots,\mathfrak E_{\nu'-1}^{(1)}$ auf, so erhalten wir Ausdrücke von der Form:
\[
L_{i,k}^{(1)}=\xi_{i,0}L_{0,k}+\xi_{i,1}L_{1,k}+\cdots+\xi_{i,\nu'-1}L_{\nu'-1,k}.
\]
Ist also $T_1$ der Regulator des Systems der $\mathfrak E_i^{(1)}$, so ergiebt sich nach dem Multiplicationssatze der Determinanten, wenn mit $C$ der absolute Werth der Determinante der ganzzahligen Exponenten $\xi_{i,k}$, also eine ganze rationale positive Zahl bezeichnet wird:
\begin{equation}
T_1=CT_0.
\tag{13}
\end{equation}

Es kann nicht bewiesen werden und trifft für die höheren Werthe von $m$ wahrscheinlich auch nicht zu, dass $\tau_0,\tau_1,\ldots,\tau_{\nu'-1}$ ein Fundamentalsystem bilden\footnote{Es ist hier, wie bei allen Fragen, die die Einheiten betreffen, sehr schwierig, zu bestimmten numerischen Resultaten zu gelangen. Nach den Resultaten, die in dem reichhaltigen Werke von Reuschle, „Tafeln complexer Primzahlen, welche aus Wurzeln der Einheiten gebildet sind“ (Berlin 1875), enthalten sind, bilden für $m=16$ die $\tau$ noch ein Fundamentalsystem. Denn hier ist die Grundzahl des Körpers $H_m$ gleich 21, der Grad des Körpers gleich 4, und daher giebt es nach der Anmerkung zu \S~156 in jeder Idealclasse ein Ideal, dessen Norm kleiner als 6 ist. Nach den Tafeln von Reuschle kann aber jede natürliche Primzahl unter 100 im Körper $H_m$ in wirklich existirende Primfactoren zerlegt werden. Folglich gehen gewiss in allen Zahlen unter 6 nur Hauptideale auf, und folglich giebt es hier nur die eine Classe, die Hauptclasse. Da, wie wir gesehen haben, auch der erste Classenzahlfactor $=1$ ist, so ist der Kreistheilungskörper $\Omega_{16}$ einclassig. Es gelten in ihm dieselben Gesetze der Primzahlen, wie im Körper der rationalen Zahlen. Zweifelhaft ist dies für den Körper $\Omega_{32}$ und sicher nicht mehr zutreffend im Körper $\Omega_{64}$, in dem die Classenzahl ein Vielfaches von 17 ist.}.

Für unseren Zweck ist es aber von Wichtigkeit, dass sich diese Einheiten in Bezug auf den Nenner 2 in den Exponenten
\srcpage{732}
gewissermaassen wie ein Fundamentalsystem verhalten, was durch folgenden Satz bestimmter ausgedrückt wird:

\begin{center}\textbf{2. Wenn eine Normaleinheit des Körpers $H_m$ von der Form
\[
\mathfrak E(r)=\tau_0^{e_0}\tau_1^{e_1}\cdots\tau_{\nu'-1}^{e_{\nu'-1}},
\]
in der die Exponenten $e_0,e_1,\ldots,e_{\nu'-1}$ ganze rationale Zahlen sind, mit allen ihren conjugirten Werthen einerlei Vorzeichen hat, so müssen die sämmtlichen $e_0,e_1,\ldots,e_{\nu'-1}$ gerade Zahlen sein, und $\mathfrak E(r)$ ist das Quadrat einer Normaleinheit.}\end{center}

Den Beweis können wir wieder durch vollständige Induction führen, wollen aber zunächst den Ausdruck für $\tau_j$ etwas umformen.

Die conjugirten Werthe zu der Einheit $\mathfrak E(r)$ erhalten wir in der Form
\[
\mathfrak E(r_\beta)=\tau_\beta^{e_0}\tau_{\beta+1}^{e_1}\cdots\tau_{\beta+\nu'-1}^{e_{\nu'-1}},
\]
und diese Ausdrücke sollen also für alle Werthe von $\beta$ dasselbe Vorzeichen haben.

Es ist nach \S~199, (5)
\[
\tau_j=\tg\frac{5^j\pi}{m}=-\frac12\frac{\sin\left(\frac{2\pi}{m}5^j\right)}{\left[\cos\left(\frac{5^j\pi}{m}\right)\right]^2}.
\]
Hierin setzen wir zur Abkürzung
\begin{equation}
\sigma_j=\sin\left(\frac{2\pi}{m}5^j\right),
\tag{14}
\end{equation}
so dass $\sigma_j$ immer dasselbe Vorzeichen hat, wie $\tau_j$.

Daraus bilden wir das Product
\begin{equation}
S_\beta=\sigma_\beta^{e_0}\sigma_{\beta+1}^{e_1}\cdots\sigma_{\beta+\nu'-1}^{e_{\nu'-1}},
\tag{15}
\end{equation}
und unsere Voraussetzung ist also die, dass $S_\beta$ für alle Werthe von $\beta$ dasselbe Vorzeichen hat.

Es soll nachgewiesen werden, dass die ganzen Zahlen $e_0,e_1,\ldots,e_{\nu'-1}$ alle gerade sein müssen.

Der Satz ist richtig für $m=8$; denn in diesem Falle ist
\[
\sigma_0=\sin\frac{2\pi}{8},\qquad \sigma_1=\sin\frac{10\pi}{8}=-\sin\frac{6\pi}{8},
\]
\srcpage{733}
also $\sigma_0$ positiv, $\sigma_1$ negativ. Hier ist nun $S_\beta=\sigma_\beta^{e_0}$, und es kann also $S_0$ und $S_1$ nur dann gleiches Vorzeichen haben, wenn $e_0$ gerade ist.

Hiernach wenden wir die vollständige Induction an. Wir setzen zur Vereinfachung $\nu'=2\nu''$ und erinnern an die Bezeichnung
\begin{equation}
m=2^\kappa,\qquad \mu=\frac12m,\qquad \nu=\frac12\mu,
\qquad \nu'=\frac12\nu,
\qquad \nu''=\frac12\nu',
\tag{16}
\end{equation}
und setzen voraus, dass unser Satz als richtig erwiesen sei, wenn $\mu$ an die Stelle von $m$ tritt.

Es ist dann
\[
5^\nu\equiv1\pmod m,\\ 5^{\nu'}\equiv1+\mu\pmod m,
\quad 5^{\nu''}\equiv1+\nu\pmod\mu,
\]
\[
5^{\nu''}\equiv1+\nu+\mu\pmod m
\]
und daraus ergiebt sich
\[
\sigma_{\beta+\nu}=-\sigma_\beta,
\]
\[
\sigma_{\beta+\nu''}=-\cos\left(\frac{2\pi}{m}5^\beta\right),
\]
\begin{equation}
\sigma_\beta\sigma_{\beta+\nu''}=-\frac12\sin\left(\frac{2\pi}{\mu}5^\beta\right)=-\frac12\sigma'_\beta,
\tag{17}
\end{equation}
\begin{equation}
\sigma'_{\beta+\nu''}=-\sigma'_\beta,
\tag{18}
\end{equation}
wenn die $\sigma'_\beta$ aus den $\sigma_\beta$ durch den Uebergang von $\kappa$ zu $\kappa-1$ hervorgehen. Nun bilden wir nach (17), (18) das Product $S_\beta S_{\beta+\nu''}$, und erhalten, wenn noch zur Abkürzung
\[
e'=e_0+e_1+\cdots+e_{\nu'-1},
\qquad e''=e_0+e_1+\cdots+e_{\nu''-1}
\]
gesetzt wird,
\[
S_\beta S_{\beta+\nu''}=(-1)^{\nu''}\left(\frac12\right)^{e'}
(\sigma'_\beta)^{e_0+e_{\nu''}}(\sigma'_{\beta+1})^{e_1+e_{\nu''+1}}\cdots
(\sigma'_{\beta+\nu''-1})^{e_{\nu''-1}+e_{\nu'-1}}.
\]

Dies ist aber ein Product von derselben Form wie $S_\beta$, in dem $\mu$ an Stelle von $m$ getreten ist, was gleichfalls mit allen seinen Conjugirten einerlei Vorzeichen hat. Es ist daher nach der gemachten Voraussetzung
\begin{equation}
e_0\equiv e_{\nu''},\ e_1\equiv e_{\nu''+1},\ldots,e_{\nu''-1}\equiv e_{\nu'-1}\pmod2.
\tag{19}
\end{equation}
Setzen wir also
\[
S'_\beta=(\sigma'_\beta)^{e_0}(\sigma'_{\beta+1})^{e_1}\cdots(\sigma'_{\beta+\nu''-1})^{e_{\nu''-1}},
\]
\[
R_\beta=\sigma_{\beta+\nu''}^{\frac12(e_0-e_{\nu''})}
\sigma_{\beta+\nu''+1}^{\frac12(e_1-e_{\nu''+1})}\cdots
\sigma_{\beta+\nu'-1}^{\frac12(e_{\nu''-1}-e_{\nu'-1})},
\]
so folgt aus (15) und (17)
\[
S_\beta R_\beta=\left(-\frac12\right)^{e''}S'_\beta,
\]
\srcpage{734}
und hieraus ergiebt sich, dass $S'_\beta$ ebenso wie $S_\beta$ mit allen seinen conjugirten Werthen dasselbe Zeichen hat. $S'_\beta$ entsteht aber aus $S_\beta$ dadurch, dass $\mu$ an Stelle von $m$ gesetzt wird, und nach unserer Voraussetzung ist daher
\[
e_0\equiv0,\ e_1\equiv0,\ldots,e_{\nu''-1}\equiv0\pmod2,
\]
und mit Hülfe von (19)
\[
e_{\nu''}\equiv0,\ e_{\nu''+1}\equiv0,\ldots,e_{\nu'-1}\equiv0\pmod2,
\]
wodurch der Beweis des Satzes 2. geführt ist.

Wenn wir nun in der Formel (12) für $\mathfrak E_0,\mathfrak E_1,\ldots,\mathfrak E_{\nu'-1}$ das System $\tau_0,\tau_1,\ldots,\tau_{\nu'-1}$ wählen und die Exponenten $\xi$ in die Form setzen:
\[
\xi_0=\frac{e_0}{e},\quad \xi_1=\frac{e_1}{e},\ldots,\quad \xi_{\nu'-1}=\frac{e_{\nu'-1}}{e},
\]
worin $e,e_0,e_1,\ldots,e_{\nu'-1}$ ganze rationale Zahlen ohne gemeinsamen Theiler sind, so erhalten wir
\[
\mathfrak E(r)^e=\pm\tau_0^{e_0}\tau_1^{e_1}\cdots\tau_{\nu'-1}^{e_{\nu'-1}},
\]
und es folgt aus unserem Satze, dass $e$ ungerade sein muss; denn wäre es gerade, so wäre $\mathfrak E(r)^e$ ein Quadrat einer reellen Grösse, und folglich hätte $\pm\mathfrak E(r)^e$ mit seinen Conjugirten dasselbe Vorzeichen. Dann aber müssten nach unserem Satze alle $e_0,e_1,\ldots,e_{\nu'-1}$ durch 2 theilbar sein, was wieder gegen die Voraussetzung ist, dass die Exponenten $e,e_0,\ldots,e_{\nu'-1}$ keinen gemeinsamen Theiler haben sollen. Wir haben daher den Satz:

\begin{center}\textbf{3. Ist $\mathfrak E(r)$ irgend eine Normaleinheit des Körpers $H_m$, so lassen sich die ungerade Zahl $e$ und die ganzen Zahlen $e_0,e_1,\ldots,e_{\nu'-1}$ so bestimmen, dass}
\[
\mathfrak E(r_\beta)^e=\pm\tau_\beta^{e_0}\tau_{\beta+1}^{e_1}\cdots\tau_{\beta+\nu'-1}^{e_{\nu'-1}}
\tag{20}
\]
\textbf{wird.}\end{center}

Dies Ergebniss führt zu einigen wichtigen Consequenzen:

\begin{center}\textbf{4. Ist $T$ der Regulator des Systems der $\tau_j$, und $T_0$ der Regulator eines Fundamentalsystems normaler Einheiten in $H_m$, d. h. der Minimalwerth, den der Regulator irgend eines Systems unabhängiger Normaleinheiten annehmen kann, so ist}
\[
T=CT_0,\\ C\equiv1\pmod2,
\tag{21}
\]
\textbf{worin $C$ eine ungerade natürliche Zahl bedeutet.}\end{center}

\srcpage{735}
Dass nämlich $T:T_0=C$ eine ganze Zahl ist, folgt aus der Formel (13).

Wenn wir aber den Satz 3. auf alle Elemente des Fundamentalsystems $\mathfrak E_i(r)$ anwenden, so ergiebt sich ein System ganzer Zahlen $e_{i,k}$ und eine ungerade Zahl $e$, so dass
\[
\mathfrak E_i(r)^e=\pm\tau_0^{e_{i,0}}\tau_1^{e_{i,1}}\cdots\tau_{\nu'-1}^{e_{i,\nu'-1}}.
\]
Nehmen wir hiervon die Logarithmen, und bilden dann den Regulator $T_0$, so ergiebt sich
\[
e^{\nu'}T_0=\pm\sum\pm e_{0,0}e_{1,1}\cdots e_{\nu'-1,\nu'-1}T,
\]
folglich
\[
e^{\nu'}=\pm C\sum\pm e_{0,0}e_{1,1}\cdots e_{\nu'-1,\nu'-1},
\]
und da hierin die linke Seite ungerade ist, so kann $C$ nicht gerade sein, w. z. b. w.

Wenn alle conjugirten Werthe der Normaleinheit $\mathfrak E(r_\beta)$ einerlei Zeichen haben, so gilt dasselbe auch von $\mathfrak E(r_\beta)^e$, und folglich müssen in diesem Falle in der Formel (20) die Exponenten $e_0,e_1,\ldots,e_{\nu'-1}$ nach 2. gerade Zahlen sein. Schreibt man dann die Formel (20) so:
\[
\mathfrak E(r)=\pm\tau_0^{e_0}\tau_1^{e_1}\cdots\tau_{\nu'-1}^{e_{\nu'-1}}\mathfrak E(r)^{1-e},
\]
so erhält man daraus den Satz:

\begin{center}\textbf{5. Eine Normaleinheit des Körpers $H_m$, die mit ihren Conjugirten gleiches Vorzeichen hat, ist, vom Vorzeichen abgesehen, ein Quadrat einer Normaleinheit.}\end{center}

\section*{\S\,201. Fundamentalsystem von Einheiten des Körpers $H_m$.}

Es bleibt noch übrig, den Nenner $D$ in der Formel für $B$ [\S~197, (6)] zu bestimmen, der durch die Gleichung
\[
E=E'D
\]
definirt war, worin $E$ und $E'$ die Regulatoren der Körper $H_m,H_\mu$ bedeuteten.

Hier handelt es sich also nicht mehr um die Normaleinheiten, sondern um die Einheiten des Körpers $H_m$ überhaupt.

\srcpage{736}
Wir nehmen ein Fundamentalsystem von Einheiten im Körper $H_m$ an:
\begin{equation}
\varepsilon_1(r),\ \varepsilon_2(r),\ldots,\varepsilon_{\nu-1}(r),
\tag{1}
\end{equation}
und die conjugirten Logarithmen dieses Systems seien:
\begin{equation}
l_{1,k},\ l_{2,k},\ldots,l_{\nu-1,k},\qquad k=0,1,2,\ldots,\nu-1.
\tag{2}
\end{equation}
Daneben betrachten wir ein Fundamentalsystem im Körper $H_\mu$
\begin{equation}
\varepsilon'_1,\ \varepsilon'_2,\ldots,\varepsilon'_{\nu'-1}
\tag{3}
\end{equation}
mit den conjugirten Logarithmen
\begin{equation}
l'_{1,k},\ l'_{2,k},\ldots,l'_{\nu'-1,k},\qquad k=0,1,2,\ldots,\nu'-1.
\tag{4}
\end{equation}

Bei der Bildung der Regulatoren wird einer der conjugirten Werthe, etwa $k=0$, weggelassen (vgl. \S~186), und wir erhalten daher
\begin{equation}
E=\pm\sum\pm l_{1,1}l_{2,2}\cdots l_{\nu-1,\nu-1},
\qquad
E'=\pm\sum\pm l'_{1,1}l'_{2,2}\cdots l'_{\nu'-1,\nu'-1}.
\tag{5}
\end{equation}

Nun fügen wir zu dem Systeme (3) ein Fundamentalsystem von Normaleinheiten des Körpers $H_m$
\begin{equation}
\mathfrak E_0,\mathfrak E_1,\ldots,\mathfrak E_{\nu'-1},
\tag{6}
\end{equation}
und wir behaupten, dass das System
\begin{equation}
\varepsilon'_1,\varepsilon'_2,
\ldots,\varepsilon'_{\nu'-1},\mathfrak E_0,\mathfrak E_1,
\ldots,\mathfrak E_{\nu'-1}
\tag{7}
\end{equation}
ein unabhängiges System von Einheiten des Körpers $H_m$ ist (wenn auch nicht ein Fundamentalsystem).

Bezeichnen wir mit
\begin{equation}
L_{0,k},L_{1,k},\ldots,L_{\nu'-1,k},\qquad k=0,1,\ldots,\nu'-1
\tag{8}
\end{equation}
die conjugirten Logarithmen des Systems (6), so ist nach dem vorigen Paragraphen
\begin{equation}
\pm\sum\pm L_{0,0}L_{1,1}\cdots L_{\nu'-1,\nu'-1}=T_0
\tag{9}
\end{equation}
der Regulator dieses Systems, und wir haben mit Rücksicht auf die Definition der Normaleinheiten [\S~200, (9)]:
\begin{equation}
L_{i,k+\nu'}=-L_{i,k},
\tag{10}
\end{equation}
und weil die $\varepsilon'_i$ als Zahlen des Körpers $H_\mu$ durch die Substitution $(r,-r)$ ungeändert bleiben:
\begin{equation}
l'_{i,k+\nu'}=l'_{i,k}.
\tag{11}
\end{equation}

\srcpage{737}
Hiernach ergiebt sich der Regulator $R$ des Systems (7) aus der Determinante:
\[
\small
\left|\begin{array}{cccc|cccc}
l'_{1,1}&l'_{2,1}&\cdots&l'_{\nu'-1,1}&L_{0,1}&L_{1,1}&\cdots&L_{\nu'-1,1}\\
\cdots&\cdots&&\cdots&\cdots&\cdots&&\cdots\\
l'_{1,\nu'-1}&l'_{2,\nu'-1}&\cdots&l'_{\nu'-1,\nu'-1}&L_{0,\nu'-1}&L_{1,\nu'-1}&\cdots&L_{\nu'-1,\nu'-1}\\
l'_{1,0}&l'_{2,0}&\cdots&l'_{\nu'-1,0}&-L_{0,0}&-L_{1,0}&\cdots&-L_{\nu'-1,0}\\
l'_{1,1}&l'_{2,1}&\cdots&l'_{\nu'-1,1}&-L_{0,1}&-L_{1,1}&\cdots&-L_{\nu'-1,1}\\
\cdots&\cdots&&\cdots&\cdots&\cdots&&\cdots\\
l'_{1,\nu'-1}&l'_{2,\nu'-1}&\cdots&l'_{\nu'-1,\nu'-1}&-L_{0,\nu'-1}&-L_{1,\nu'-1}&\cdots&-L_{\nu'-1,\nu'-1}
\end{array}\right|.
\]
Diese Determinante hat $\nu-1$ Zeilen; wir addiren die erste zur $(\nu'+1)$ten, die zweite zur $(\nu'+2)$ten u.s.f., die $(\nu'-1)$te zur letzten. Dadurch heben sich in den letzten $(\nu'-1)$ Zeilen die $L$ heraus und die $l'$ bekommen den Factor 2 [nach (10), (11)], und hieraus ergiebt sich nach einem Determinantensatze (Bd. I, \S~23, V.):
\begin{equation}
R=2^{\nu'-1}E'T_0.
\tag{12}
\end{equation}

Da hiernach $R$ nicht verschwindet und daher in (7) ein System von $\nu-1$ unabhängigen Einheiten des Körpers $H_m$ vorliegt, so können wir [nach \S~186, (5)] die rationalen (ganzen oder gebrochenen) Zahlen $m_{i,k},M_{i,k}$ aus den Gleichungen
\begin{equation}
2l_{i,k}=m_{1,i}l'_{1,k}+\cdots+m_{\nu'-1,i}l'_{\nu'-1,k}
+M_{0,i}L_{0,k}+M_{1,i}L_{1,k}+\cdots+M_{\nu'-1,i}L_{\nu'-1,k}
\tag{13}
\end{equation}
bestimmen. Es lässt sich aber leicht durch die Relationen (10), (11) nachweisen, dass die Zahlen $m_{i,k},M_{i,k}$ hier \textbf{ganze} sein müssen. Denn es ist
\[
l_{i,k}+l_{i,k+\nu'}=\log|\varepsilon_i(r_k)\varepsilon_i(-r_k)|
=m_{1,i}l'_{1,k}+\cdots+m_{\nu'-1,i}l'_{\nu'-1,k},
\]
\[
l_{i,k}-l_{i,k+\nu'}=\log|\varepsilon_i(r_k)\varepsilon_i(-r_k)^{-1}|
=M_{0,i}L_{0,k}+M_{1,i}L_{1,k}+\cdots+M_{\nu'-1,i}L_{\nu'-1,k}.
\]
Die Einheiten $\varepsilon_i(r)\varepsilon_i(-r)$ gehören aber dem Körper $H_\mu$ an, und da die $\varepsilon'_i$ nach Voraussetzung ein Fundamentalsystem dieses Körpers sind, so müssen die $m_{i,k}$ nach \S~187 ganze rationale Zahlen sein.

Die Einheiten $\varepsilon_i(r)\varepsilon_i(-r)^{-1}$ sind Normaleinheiten des Körpers $H_m$, und weil die $\mathfrak E_i$ ein Fundamentalsystem von Normaleinheiten sein sollten, so sind nach \S~200 die $M_{i,k}$ ganze Zahlen.

\srcpage{738}
Zur besseren Uebersicht setzen wir die Gleichungen (13) noch einmal ohne den Index $k$, der die conjugirten Werthe von einander unterscheidet, hierher:
\begin{equation}
2l_i=m_{1,i}l'_1+m_{2,i}l'_2+\cdots+m_{\nu'-1,i}l'_{\nu'-1}
+M_{0,i}L_0+M_{1,i}L_1+\cdots+M_{\nu'-1,i}L_{\nu'-1}.
\tag{14}
\end{equation}
Bilden wir nun nach (13) den Regulator $E$ und bezeichnen mit $M$ den absoluten Werth der Determinante der $\nu-1$ Reihen der Zahlen $m_{k,i},M_{k,i}$:
\[
M=\pm\sum\pm m_{1,1}\cdots m_{\nu'-1,\nu'-1}M_{0,0}M_{1,1}\cdots M_{\nu'-1,\nu'-1},
\]
so folgt:
\begin{equation}
2^{\nu-1}E=MR,
\tag{15}
\end{equation}
und daraus mit Rücksicht auf (12):
\begin{equation}
E=2^{-\nu'}ME'T_0.
\tag{16}
\end{equation}

Es lässt sich nun weiter nachweisen, dass $M$ eine Potenz von 2 und durch $2^{\nu'}$ theilbar ist.

Da nämlich das System (1) als Fundamentalsystem in $H_m$ vorausgesetzt war, so giebt es ganze rationale Zahlen $n_{k,i},N_{k,i}$, die den Gleichungen genügen:
\begin{equation}
\begin{aligned}
l'_i&=n_{1,i}l_1+n_{2,i}l_2+\cdots+n_{\nu-1,i}l_{\nu-1},&& i=1,2,\ldots,\nu'-1,\\
L_i&=N_{1,i}l_1+N_{2,i}l_2+\cdots+N_{\nu-1,i}l_{\nu-1},&& i=0,1,2,\ldots,\nu'-1.
\end{aligned}
\tag{17}
\end{equation}
Bilden wir daraus die Determinante $R$ und bezeichnen mit $N$ den absoluten Werth der Determinante der $n_{k,i},N_{k,i}$, so folgt:
\[
R=NE,
\]
und daraus nach (15):
\[
MN=2^{\nu-1};
\]
daraus folgt, dass jede der beiden ganzen Zahlen $M,N$ eine Potenz von 2 sein muss.

Um zu entscheiden, durch welche Potenz von 2 die Determinante $M$ theilbar ist, bestimmen wir zunächst ein System ganzer rationaler Zahlen $a_i$ ohne gemeinschaftlichen Theiler, das den linearen Gleichungen
\begin{equation}
\sum_{i=1}^{\nu-1}a_im_{k,i}=0\qquad(k=1,2,\ldots,\nu'-1)
\tag{18}
\end{equation}
genügt. Setzen wir dann
\[
\sum_{i=1}^{\nu-1}a_iM_{k,i}=\xi_k,
\]
\srcpage{739}
so erhalten wir aus den Gleichungen (13)
\[
2\sum_{i=1}^{\nu-1}a_il_{i,k}=\xi_0L_{0,k}+\xi_1L_{1,k}+\cdots+\xi_{\nu'-1}L_{\nu'-1,k}.
\]
Daraus ergiebt sich nach (10)
\[
\sum_{i=1}^{\nu-1}a_il_{i,k+\nu'}=-\sum_{i=1}^{\nu-1}a_il_{i,k},
\]
und dies bedeutet, dass
\[
\varepsilon_1^{a_1}\varepsilon_2^{a_2}\cdots\varepsilon_{\nu-1}^{a_{\nu-1}}
\]
eine Normaleinheit des Körpers $H_m$ ist. Daraus folgt dann weiter, weil $\mathfrak E_0,\mathfrak E_1,\ldots,\mathfrak E_{\nu'-1}$ ein Fundamentalsystem von Normaleinheiten ist, dass die Zahlen $\xi_0,\xi_1,\ldots,\xi_{\nu'-1}$ gerade sein müssen. Daraus ergiebt sich:

\begin{center}\textbf{1. Das System der Gleichungen (18) hat die Congruenzen
\[
\sum_{i=1}^{\nu-1}a_iM_{k,i}\equiv0\pmod2\qquad(k=0,1,\ldots,\nu'-1)
\tag{19}
\]
zur Folge.}\end{center}

In den Gleichungen (18), deren Anzahl nur $\nu'-1$ beträgt, sind aber $\nu-1$ Unbekannte $a_i$ enthalten. Daher giebt es mehrere Systeme von einander unabhängiger Lösungen, was wir etwas genauer dahin präcisiren können:

\begin{center}\textbf{2. Es giebt $\nu'$ ganzzahlige Lösungen des Systems (18):}
\[
\begin{array}{cccc}
a_{1,1}&a_{2,1}&\cdots&a_{\nu-1,1}\\
a_{1,2}&a_{2,2}&\cdots&a_{\nu-1,2}\\
\cdots&\cdots&&\cdots\\
a_{1,\nu'}&a_{2,\nu'}&\cdots&a_{\nu-1,\nu'}
\end{array}
\tag{20}
\]
\textbf{von der Art, dass aus der Matrix (20) eine Determinante von $\nu'$ Reihen gebildet werden kann, die eine ungerade Zahl ist.}\end{center}

Denn es giebt zunächst eine Lösung von (18):
\[
a_{1,1},a_{2,1},\ldots,a_{\nu-1,1},
\]
in der eine ungerade Zahl vorkommt. Nehmen wir an, was gestattet ist, da hier auf die Anordnung der Indices nichts ankommt,
\srcpage{740}
es sei $a_{1,1}$ diese ungerade Zahl, dann bestimmen wir eine zweite Lösung:
\[
0,a_{2,2},a_{3,2},\ldots,a_{\nu-1,2},
\]
in der wieder eine ungerade Zahl, etwa $a_{2,2}$, vorkommt. Dann bestimmen wir eine dritte Lösung:
\[
0,0,a_{3,3},\ldots,
\]
und fahren mit Nullsetzen so lange fort, als die Anzahl der übrig bleibenden Unbekannten noch grösser ist, als die Anzahl der Gleichungen.

Im letzten Systeme werden also $\nu'-1$ Nullen vorkommen, damit $\nu'$ Unbekannte übrigbleiben.

Diese Zahlenreihen $a_{i,k}$ können wir für das System (20) setzen, denn darin reducirt sich die Determinante
\begin{equation}
\pm\sum\pm a_{1,1}a_{2,2}\cdots a_{\nu',\nu'}
\tag{21}
\end{equation}
auf das Diagonalglied $a_{1,1}a_{2,2}\cdots a_{\nu',\nu'}$, dessen Factoren alle ungerade sind.

Ist nun die Forderung des Satzes 2. erfüllt, so lässt sich die Matrix (20) durch Hinzufügung von $\nu'-1$ Zeilen zu einer Determinante von $\nu-1$ Reihen ergänzen, die gleichfalls einen ungeraden ganzzahligen Werth hat:
\begin{equation}
\alpha=\pm\sum\pm a_{1,1}a_{2,2}\cdots a_{\nu',\nu'}a_{\nu'+1,\nu'+1}\cdots a_{\nu-1,\nu-1}.
\tag{22}
\end{equation}
Man braucht nur, wenn die Determinante (21) als ungerade vorausgesetzt wird, in den hinzugefügten Zeilen für die Elemente lauter Nullen zu setzen, ausgenommen die in der Diagonalreihe stehenden, $a_{\nu'+1,\nu'+1},\ldots,a_{\nu-1,\nu-1}$, die gleich 1 genommen werden können. Dann erhält $\alpha$ denselben Werth wie die Determinante (21).

Wenn wir aber jetzt nach dem Multiplicationssatze der Determinanten das Product $\alpha M$ bilden, so ergiebt sich eine Determinante, in der die $\nu-1$ Summen
\[
\sum_{i=1}^{\nu-1}a_i m_{k,i},
\qquad
\sum_{i=1}^{\nu-1}a_i M_{k,i},
\]
\[
k=1,2,\ldots,\nu'-1,
\qquad
k=0,1,2,\ldots,\nu'-1
\]
für jeden Werth $s=1,2,\ldots,\nu'$ eine Reihe bilden, und in der also $\nu'$ Reihen vorkommen, deren Elemente alle durch 2 theilbar sind.

\srcpage{741}
Demnach ist $\alpha M$ und folglich auch $M$ durch $2^{\nu'}$ theilbar, und wenn wir
\begin{equation}
M=2^{\nu'+\theta},\qquad N=2^{\nu'-\theta-1}
\tag{23}
\end{equation}
setzen, so ist $\theta$ eine nicht negative ganze rationale Zahl.

Setzen wir dies in (16) ein, so folgt:
\[
E=2^\theta E'T_0.
\]
Nun war in \S~197, (6)
\[
E=E'D
\]
gesetzt, und es ergiebt sich also:
\[
D=2^\theta T_0.
\]
Hieraus folgt für den Classenzahlfactor $B$ nach \S~199, (15) und \S~200, (21):
\begin{equation}
B=2^{-\theta}\frac{T}{T_0}=2^{-\theta}C,
\tag{24}
\end{equation}
worin $C$ eine ungerade ganze Zahl ist.

Demnach ergiebt sich nach \S~197, (4) für den zweiten Factor der Classenzahl $h_1$, der, wie wir gesehen haben, eine ganze Zahl ist:
\begin{equation}
h_1=2^{-\theta}h'_1C.
\tag{25}
\end{equation}
Nehmen wir nun als bewiesen an, dass $h'_1$ ungerade ist, so folgt, da auch $C$ ungerade ist, dass $\theta=0$ und $h_1$ ungerade ist, und beides ist also damit durch vollständige Induction bewiesen.

Damit werden die abgeleiteten Resultate folgendermaassen vereinfacht:

\begin{center}\textbf{3. Es ist}
\[
M=2^{\nu'},\qquad D=T_0,
\qquad B=\frac{T}{T_0};
\tag{26}
\]
\textbf{der Classenzahlfactor $B$ ist gleich dem Verhältniss des Regulators $T$ der Einheiten $\tau$ zum Regulator $T_0$ eines Fundamentalsystems von Normaleinheiten, und ist eine ungerade ganze Zahl.}\end{center}

Bezeichnen wir die zu $m=2^\kappa$ gehörige Zahl $B$ mit $B_\kappa$, so ist die Classenzahl $h_1$, wie durch vollständige Induction folgt,
\begin{equation}
h_1=B_4B_5\cdots B_\kappa,
\tag{27}
\end{equation}
und ist also auch eine ungerade Zahl. Damit ist dann, mit Rücksicht auf \S~198, 1., das Haupttheorem bewiesen:

\begin{center}\textbf{A. Die Classenzahl im Körper $\Omega_m$ ist, wenn $m$ eine Potenz von 2 ist, ungerade.}\end{center}

\section*{\S\,202. Positive Einheiten.}

\srcpage{742}
Die zuletzt gefundenen Resultate zeigen, dass das System unabhängiger Einheiten \S~201, (7)
\[
\varepsilon'_1,\varepsilon'_2,\ldots,\varepsilon'_{\nu'-1},\mathfrak E_0,\mathfrak E_1,\ldots,\mathfrak E_{\nu'-1}
\]
im Körper $H_m$ kein Fundamentalsystem bildet. Denn sonst müssten in den Formeln \S~201, (13) die Zahlen $m_{k,i},M_{k,i}$ sämmtlich durch 2 theilbar sein, und es wäre die Determinante $M$, deren Werth wir $=2^{\nu'}$ gefunden haben, durch $2^{\nu-1}$ theilbar.

Wir betrachten nun an Stelle des Systems der Gleichungen (18), \S~201 die folgenden Congruenzen:
\begin{equation}
\begin{aligned}
\sum_{i=1}^{\nu-1}b_i m_{1,i}&\equiv1\pmod2,\\
\sum_{i=1}^{\nu-1}b_i m_{2,i}&\equiv0\pmod2,\\
&\ \vdots\\
\sum_{i=1}^{\nu-1}b_i m_{\nu'-1,i}&\equiv0\pmod2,
\end{aligned}
\tag{1}
\end{equation}
und beweisen, dass es möglich ist, ihnen durch ganzzahlige Werthe der $b_i$ zu genügen.

Wäre es nicht möglich, so müsste sich aus den $\nu'-2$ Congruenzen
\[
\sum_{i=1}^{\nu-1}b_i m_{k,i}\equiv0\pmod2\qquad k=2,3,\ldots,\nu'-1
\]
die letzte
\[
\sum_{i=1}^{\nu-1}b_i m_{1,i}\equiv0\pmod2
\]
als Folgerung ergeben, und dann würde, wie im vorigen Paragraphen, weiter folgen:
\[
\sum_{i=1}^{\nu-1}b_i M_{k,i}\equiv0\pmod2\qquad k=0,1,\ldots,\nu'-1.
\]

\srcpage{743}
Man könnte dann an Stelle der Matrix \S~201, (20) eine Matrix der $b_{k,i}$ von $\nu'+1$ Reihen setzen, und es würde sich auf dem im vorigen Paragraphen eingeschlagenen Wege schliessen lassen, dass $M$ durch $2^{\nu'+1}$ theilbar sein müsste, was nicht der Fall ist.

Wenden wir aber das System der Multiplicatoren $b_i$, das den Bedingungen (1) genügt, auf die Gleichungen (14) des vorigen Paragraphen an, indem wir mit $b_i$ multipliciren und dann summiren, so ergiebt sich, wenn wir Vielfache von 2 weglassen und dies auch hier (wo irrationale Zahlen, nämlich die Logarithmen, auftreten) durch das Zeichen der Congruenz ausdrücken:
\begin{equation}
l'_1\equiv L_0\sum M_{0,i}b_i+L_1\sum M_{1,i}b_i+\cdots+L_{\nu'-1}\sum M_{\nu'-1,i}b_i\pmod2.
\tag{2}
\end{equation}
Wollen wir die Congruenz in eine Gleichung verwandeln, so kommt eine Summe von Grössen
\[
l_1,l_2,\ldots,l_{\nu-1},\quad l'_1,l'_2,\ldots,l'_{\nu'-1}
\]
mit geraden ganzen Zahlen als Coefficienten hinzu.

Dieselbe Betrachtung lässt sich aber auf $l'_2,l'_3,\ldots,l'_{\nu'-1}$ anwenden, und danach kann man, wenn man von den Logarithmen wieder zu den Zahlen übergeht, die Formel (2) so in Worte fassen:

\begin{center}\textbf{1. Jede Einheit des Körpers $H_\mu$ ist als Product einer Normaleinheit des Körpers $H_m$ und eines Quadrates einer Einheit in $H_m$ darstellbar.}\end{center}

Wenn wir darauf den Satz 2., \S~200 anwenden, so ergiebt sich, dass jede Einheit $\mathfrak E(r)$ des Körpers $H_\mu$, die mit allen ihren Conjugirten dasselbe Zeichen hat, vom Vorzeichen abgesehen, das Quadrat einer Einheit in $H_m$ sein muss. Ist aber
\[
\pm\mathfrak E(r)=\Delta(r)^2,
\]
worin $\mathfrak E(r)$ eine Einheit in $H_\mu$ und $\Delta(r)$ eine Einheit in $H_m$ ist, so muss auch $\Delta(r)$ in $H_\mu$ enthalten sein.

Denn setzen wir
\[
\Delta(r)=\Delta_1(r^2)+r\Delta_2(r^2),
\]
so kann das Quadrat davon nur dann in $H_\mu$ enthalten sein, also durch die Vorzeichenänderung von $r$ ungeändert bleiben, wenn entweder $\Delta_1$ oder $\Delta_2=0$ ist. Es kann aber nicht $\Delta(r)=r\Delta_2(r^2)$ sein, denn sonst wäre $\Delta_2(r^2)$ eine Einheit in $\Omega_\mu$ und müsste
\srcpage{744}
nach \S~178, 1. das Product einer reellen Einheit und einer Einheitswurzel $r^{2\lambda}$ sein, und da auch $\Delta(r)$ reell ist, so müsste $r^{2\lambda+1}$ reell sein, was nicht möglich ist. Demnach ist $\Delta(r)$ in $H_\mu$ enthalten.

Da nun $\mu$ ebenso wie $m$ jede beliebige Potenz von 2 sein kann, so ergiebt sich hieraus das zweite Haupttheorem:

\begin{center}\textbf{B. Eine Einheit des Körpers $H_m$, die mit allen ihren Conjugirten positiv ist, ist das Quadrat einer Einheit desselben Körpers.}\end{center}

Von den beiden Theoremen A. und B. haben wir aber im \S~184 den Beweis des Hauptsatzes von den Abel'schen Zahlkörpern (\S~179, I.) abhängig gemacht.

\end{document}
